%% ─────────────────────────────────────────────────────────────────────────────
%% CV — Talvin Ackbaraly
%% Stack : LaTeX (pdflatex) — 1 Page Strict Architecture
%% ─────────────────────────────────────────────────────────────────────────────
\documentclass[10pt, a4paper]{article}

%% ── Packages ─────────────────────────────────────────────────────────────────
\usepackage[T1]{fontenc}
\usepackage[utf8]{inputenc}
\usepackage[french]{babel}
\usepackage[
  top=1.4cm, bottom=1.4cm,
  left=1.6cm, right=1.6cm
]{geometry}
\usepackage{microtype}

% Police moderne Sans-Serif (Helvetica / TeX Gyre Heros)
\usepackage{tgheros}
\renewcommand{\familydefault}{\sfdefault}

\usepackage{parskip}
\usepackage[explicit]{titlesec}
\usepackage{enumitem}
\usepackage{tabularx}
\usepackage{array}
\usepackage{xcolor}
\usepackage{hyperref}
\usepackage{fontawesome5}
\usepackage{tikz}

%% ── Palette de Couleurs ──────────────────────────────────────────────────────
\definecolor{accent}{HTML}{0284c7}      % Saphir Blue
\definecolor{dark}{HTML}{0f172a}        % Slate 900
\definecolor{textmain}{HTML}{334155}    % Slate 700
\definecolor{muted}{HTML}{64748b}       % Slate 500
\definecolor{subtle}{HTML}{94a3b8}      % Slate 400
\definecolor{rule}{HTML}{e2e8f0}        % Slate 200
\definecolor{tagbg}{HTML}{f1f5f9}       % Slate 100

%% ── Hyperliens ───────────────────────────────────────────────────────────────
\hypersetup{
  colorlinks = true,
  urlcolor   = accent,
  linkcolor  = accent,
  pdftitle   = {CV — Talvin Ackbaraly},
  pdfauthor  = {Talvin Ackbaraly},
}

\color{textmain}

%% ── Typographie des sections ─────────────────────────────────────────────────
\titleformat{\section}[block]
  {\Large\bfseries\color{dark}}
  {}
  {0em}
  {#1\\[-0.2\baselineskip]{\color{accent}\titlerule[1.5pt]}\vspace{4pt}}

\titlespacing*{\section}{0pt}{14pt}{0pt}

%% ── Helpers & Micro-Components ───────────────────────────────────────────────
\newcommand{\cvtag}[1]{%
  \tikz[baseline]\node[anchor=base, draw=none, fill=tagbg, rounded corners=3pt, inner xsep=5pt, inner ysep=3pt, text height=1.4ex, text depth=0.2ex, text=muted, font=\scriptsize\bfseries] {#1};%
}

\newcommand{\period}[1]{\textcolor{subtle}{\small #1}}
\newcommand{\company}[1]{\textbf{\textcolor{accent}{#1}}}

\newcommand{\projectItem}[3]{%
  \noindent
  \begin{tabularx}{\textwidth}{@{}p{1.5cm} X >{\raggedleft\arraybackslash}p{3.5cm}@{}}
    \textbf{#1} & #2 & #3
  \end{tabularx}%
  % \vspace{2pt}
}

%% ── En-tête de ligne ─────────────────────────────────────────────────────────
\newcommand{\expline}[3]{%
  \textbf{#1} \hfill \period{#3}\\
  \company{#2}%
}

%% ── Config ───────────────────────────────────────────────────────────────────
\pagestyle{empty}
\setlength{\parindent}{0pt}

\begin{document}

%% ══ EN-TÊTE ══════════════════════════════════════════════════════════════════
\begin{tabularx}{\textwidth}{@{}X r@{}}
  \begin{minipage}[t]{\linewidth}
    \vspace{0pt}
    {\Huge\bfseries\color{dark}Talvin Ackbaraly}

    \vspace{4pt}
    {\large\color{muted}Étudiant en 5e année à EPITA — Ingénierie informatique}

    \vspace{6pt}
    \tikz\node[draw=accent!30, fill=accent!8, rounded corners=4pt, inner xsep=8pt, inner ysep=4pt, text=accent, font=\footnotesize\bfseries] {\faIcon{circle}~En recherche de stage de fin d'études de 6 mois — à partir de début fév. 2027};
  \end{minipage}
  &
  \begin{tabular}[t]{@{}r@{}}
    \faIcon{globe}~\href{https://talvin-ackbaraly.com}{talvin-ackbaraly.com} \\[3.5pt]
    \faIcon{linkedin}~\href{https://www.linkedin.com/in/talvin-ackbaraly}{linkedin.com/in/talvin-ackbaraly} \\[3.5pt]
    \faIcon{envelope}~\href{mailto:talvin.ackbaraly@icloud.com}{talvin.ackbaraly@icloud.com} \\[3.5pt]
    \faIcon{phone}~\href{tel:+33769389008}{07 69 38 90 08} \\[3.5pt]
    \faIcon{github}~\href{https://github.com/talvinckb}{github.com/talvinckb}
  \end{tabular}
\end{tabularx}

\vspace{6pt}

%% ══ FORMATION ════════════════════════════════════════════════════════════════
\section{Formation}

\expline{Diplôme d'ingénieur — Spécialisation Image}{EPITA — Paris}{Sept. 2022 - 2027}\\[2pt]
Traitement \& synthèse d'images, intelligence artificielle et calcul haute performance sur GPU.

\vspace{6pt}

\expline{Échange académique}{Université du Québec à Chicoutimi (UQAC) — Canada}{Janv. - Mai 2024}

\vspace{6pt}

\textcolor{muted}{\faIcon{certificate}~\textbf{Anglais} : certifié TOEIC}

%% ══ EXPÉRIENCES PROFESSIONNELLES ═════════════════════════════════════════════
\section{Expériences professionnelles}

\textbf{Développeur Full-Stack \& DevOps (Stage)} \hfill \cvtag{Full-Stack} \cvtag{DevOps} \cvtag{CI/CD} \cvtag{AWS} \cvtag{Agile} \\[2pt]
\company{Studiocanal} \hfill \period{Sept. 2025 - Janv. 2026}\\[3pt]
Développement de modules full-stack pour la gestion des droits, contrats et distribution du catalogue de films. Mise en place et maintenance d'infrastructures cloud (CI/CD, AWS) au sein d'une équipe Agile.

\vspace{6pt}

\textbf{Consultant en informatique (Stage)} \hfill \cvtag{Django} \cvtag{Python} \cvtag{Agile SAFe}\\[2pt]
\company{Hitachi Solutions France} \hfill \period{Juin - Juil. 2022}\\[3pt]
Prototype d'application de gestion (fournisseurs, produits, contrats) pour le client Sonepar (Django, Agile SAFe).

%% ══ PROJETS ACADÉMIQUES ══════════════════════════════════════════════════════
\section{Projets académiques}

\projectItem{PFEE}{Projet de fin d'études en partenariat avec la BNF : segmentation et classification d'illustrations extraites de documents patrimoniaux numérisés. Benchmarks de modèles (YOLO, Florence-2), définition des métriques, suivi client hebdomadaire avec la BNF. \textit{(équipe de 4)}}{\cvtag{Deep Learning} \cvtag{YOLO} \cvtag{Gestion de projet}}

\projectItem{ALPR}{Reconnaissance automatique de plaques d'immatriculation. Détection classique (filtres, morphologie) et ML (86\% de précision). Benchmarks Python vs C++. Workflow GitLab complet (issues, MRs, CI). \textit{(équipe de 2)}}{\cvtag{OpenCV} \cvtag{C++} \cvtag{ML} \cvtag{GitLab CI}}

\projectItem{IRGPU}{Détection de mouvement vidéo temps réel. Portage CPU $\rightarrow$ GPU avec un gain de performance de 24$\times$ (analyse des goulots via Nsight). \textit{(équipe de 4)}}{\cvtag{CUDA} \cvtag{C++} \cvtag{Nsight}}

\projectItem{FluidSim}{Moteur de simulation de fluides temps réel basé sur SPH. Calcul physique géré via Compute Shaders et rendu via Framebuffers. \textit{(équipe de 2)}}{\cvtag{C++} \cvtag{OpenGL} \cvtag{GLSL}}

\projectItem{MedViz}{Prédiction de la Capacité Vitale Forcée à partir de scanners CT. Segmentation pulmonaire et ML (XGBoost). Conteneurisation Docker, pipeline GitLab CI. \textit{(équipe de 4)}}{\cvtag{Computer Vision} \cvtag{React} \cvtag{Docker} \cvtag{GitLab CI}}

%% ══ COMPÉTENCES ══════════════════════════════════════════════════════════════
\section{Compétences}

\begin{tabularx}{\textwidth}{@{}p{5.2cm} X@{}}

  \textbf{Vision par ordinateur \& IA}
  &
  \cvtag{Traitement d'image} \cvtag{Deep Learning} \cvtag{Machine Learning} \cvtag{OpenCV} \cvtag{PyTorch}
  \\[6pt]

  \textbf{Calcul haute performance}
  &
  \cvtag{CUDA} \cvtag{GPU Computing} \cvtag{OpenGL} \cvtag{GLSL} \cvtag{Nsight}
  \\[6pt]

  \textbf{Développement \& DevOps}
  &
  \cvtag{C++} \cvtag{Python} \cvtag{Docker} \cvtag{GitLab CI} \cvtag{AWS} \cvtag{React/Next.js} \cvtag{Django}

  \\

\end{tabularx}

\end{document}
